\section{Literatür ve Veri Setleri}
\label{CH:Literatur}

\subsection{İşaret dili tanıma yaklaşımları}

İzole işaret dili tanımada iki ana yaklaşım bulunmaktadır. Görüntü tabanlı
yöntemler video karelerini doğrudan evrişimli ya da üç boyutlu evrişimli ağlarla
işler. Bu yöntemler yüksek doğruluk verse de büyük modeller ve güçlü donanım
gerektirir; çoğunlukla sunucu tarafında çalıştırılır. İskelet tabanlı yöntemler ise
önce her karede el, yüz ve gövde noktalarını çıkarır, sınıflandırmayı bu koordinat
dizileri üzerinde yapar. Girdi boyutu çok küçük olduğu için modeller hafif kalır ve
kişinin görünüşü, arka plan ya da ışık gibi etkenlerden daha az etkilenir.

AUTSL veri setinin yayımlandığı çalışmada RGB ve derinlik verisi üzerinde CNN ile
LSTM birleşimi temel yöntem olarak sunulmuştur~\cite{AUTSL2020}. Aynı veri setiyle
düzenlenen ChaLearn LAP 2021 yarışmasında işaretçiden bağımsız değerlendirme
yapılmış, en iyi sonuçlar birden fazla modaliteyi (RGB, derinlik, iskelet) birleştiren
yöntemlerle alınmıştır~\cite{ChaLearn2021}. Bu yöntemlerden SAM-SLR, iskelet
noktalarını çizge evrişimli ağlarla işleyerek diğer modalitelerle
birleştirmiştir~\cite{SAMSLR2021}. Yarışmanın ilk sıraları yüzde 98 civarında test
doğruluğuna ulaşmış, ancak bu sonuçlar büyük model toplulukları ve tam video girdisiyle
elde edilmiştir; telefonda gerçek zamanlı çalışma hedeflenmemiştir.

Mobil cihaz hedefli çalışmalarda Google'ın Kaggle üzerinde düzenlediği Amerikan
İşaret Dili yarışması öne çıkmaktadır~\cite{KaggleASL2023}. Bu yarışmada girdiler
MediaPipe ile çıkarılmış nokta dizileridir ve modellerin TensorFlow Lite biçiminde,
sınırlı boyut ve sürede çalışması şart koşulmuştur. Öne çıkan çözümlerde nokta
dizileri üzerinde 1B evrişim ve küçük Transformer katmanları kullanılmıştır. Bu
projede karşılaştırılan ikinci mimari bu yaklaşımdan esinlenmiştir.

TİD üzerine iskelet tabanlı çalışmalar da son yıllarda artmaktadır. TSLFormer,
AUTSL üzerinde MediaPipe noktalarını kullanan hafif bir Transformer
önermiştir~\cite{TSLFormer2025}. Torun ve Karacı, MediaPipe Holistic noktalarıyla 41
TİD ifadesi üzerinde GRU, LSTM ve bunların birleşimlerini
karşılaştırmıştır~\cite{Torun2024}.

\subsection{Nokta çıkarımı: MediaPipe}

MediaPipe, Google'ın cihaz üzerinde çalışan algılama hatları için geliştirdiği bir
çerçevedir~\cite{MediaPipe2019}. El modeli önce avuç içini bulur, ardından her el için
21 nokta üretir~\cite{MediaPipeHands2020}. Gövde modeli (BlazePose) 33 nokta
üretir~\cite{BlazePose2020}. Bu projede yeni \textit{Tasks} arayüzündeki
\texttt{hand\_landmarker} ve \texttt{pose\_landmarker\_lite} modelleri kullanılmıştır.
Aynı modellerin hem Python'da (eğitim verisinin çıkarımı) hem Android'de (canlı
kullanım) bulunması, eğitim ile uygulama arasındaki nokta sırası ve ölçek farklarını
ortadan kaldırmak için belirleyici olmuştur.

\subsection{Türk İşaret Dili veri setleri}

Proje başında ve sonunda yapılan taramalarda bulunan kelime düzeyindeki TİD veri
setleri Tablo~\ref{tab:verisetleri} içinde özetlenmiştir.

\begin{table}[H]
\centering
\caption{Kelime düzeyindeki Türk İşaret Dili veri setleri}
\label{tab:verisetleri}
\small
\begin{tabularx}{\textwidth}{@{}l r r r X@{}}
\toprule
\textbf{Veri seti} & \textbf{İşaret} & \textbf{İşaretçi} & \textbf{Video} & \textbf{Erişim} \\
\midrule
AUTSL~\cite{AUTSL2020} & 226 & 43 & 36.302 & Akademik kullanım; bu projede kullanıldı \\
BosphorusSign22k~\cite{Bosphorus2020} & 744 & 6 & 22.542 & Yazarlara lisans sözleşmesi gönderilerek \\
TurkSign446~\cite{TurkSign446} & 446 & 10 & 44.600 & Dosyalar kısıtlı, erişim talebiyle; yalnızca MediaPipe noktaları \\
\bottomrule
\end{tabularx}
\end{table}

AUTSL, işaretçi sayısının fazlalığı ve resmî bölmelerinin işaretçiden bağımsız olması
nedeniyle ana veri seti olarak seçilmiştir. Altı işaretçili BosphorusSign22k ve on
katılımcılı TurkSign446 daha geniş sözlük sunmakla birlikte kişi çeşitliliği daha
düşüktür; ikisi de erişim izni gerektirdiği için bu dönemde kullanılamamıştır. Bu veri
setleriyle sözlüğün genişletilmesi gelecek çalışma olarak bırakılmıştır.

\subsection{Projenin farkı}

Literatürdeki yüksek doğruluklu çalışmaların önemli bir kısmı sunucuda çalışan büyük
modellere dayanmaktadır. Bu projenin katkısı modelin kendisinden çok, hattın uçtan uca
ve telefonda çalışmasındadır:

\begin{itemize}
  \item Girdi olarak video karesi değil iskelet koordinatları kullanılmış, model 1,6 MB
        ile sınırlı tutulmuştur. Ham görüntü işlenmediği için kişisel veri yükü azalmıştır.
  \item Nokta çıkarımı eğitimde ve telefonda aynı MediaPipe Tasks modelleriyle
        yapılmış, ön işleme Python ve Kotlin'de birebir aynı yazılmıştır. Aynı 200 örnekte
        iki taraf yüzde 99,5 oranında aynı tahmini vermektedir.
  \item Sistem tamamen çevrimdışıdır; uygulamanın internet izni bulunmamaktadır.
  \item Kullanıcının işareti doğru yapabilmesi için her kelimenin örnek hareketi
        uygulama içinde iskelet animasyonu olarak gösterilmektedir.
\end{itemize}
